\chapter{Background}

\fei{You need to partition this content into two chapters: 
Chapter 2 Background will have 2.1, 2.2, 2.4, just describing defintions and language used in your report.  Give the background for a reader to understand (not your new work).}

\fei{CHapter 3 Related work will talk about other related methods and how they work, specifically, strengths, weaknesses and how your tool is different.  Do this for each method.  This will include Sec. 2.3, and 2.4}

\fei{Section 2.6 and 2.7 are more implementation details, which you describe in your later Chapter on system implementation.}

\section{Time series}
\begin{itemize}
    \item What is time series
\end{itemize}

\section{Time series anomalies}
\begin{itemize}
    \item What is time series anomalies
    \item Anomalies types: point anomalies, sequential anomalies, ...
    \item Classification of anomalies in time series
\end{itemize}

\section{Time series anomaly detection algorithms}
\begin{itemize}
    \item Introduce some time series anomaly detection algorithms (SAND, DUMP, NormA, etc.)
\end{itemize}

\section{Time series drift}
\begin{itemize}
    \item What is time series anomalies drift
    \item Drift types: abrupt, gradual, recurring
    \item Concept drift affects the accuracy of time series anomaly detection algorithms.
\end{itemize}

\section{Time series drift detection algorithms}
\begin{itemize}
    \item Introduce some time series drift detection (DDG-DA)
\end{itemize}


\section{Web tools for plotting interactive time series data}
\begin{itemize}
    \item Tools for plotting interactive time series data
    \item Web-based tools such as JavaScript libraries (e.g., ECharts, D3.js, Plotly.js)
    \item Python plotting libraries (matplotlib.py)
    \item Why choose JavaScript libraries, compare native Python libraries and JavaScript libraries
\end{itemize}

\section{Backend framework}
\begin{itemize}
    \item Introduce Python backend framework libraries, and compare them, why choose flask.py
\end{itemize}


